\section*{第 \ref{ch:b2:boltzmann-factor} 章——玻尔兹曼因子}

\begin{solution}{exo:b2:boltzmann-factor:1}
$H = k_BT/mg = \qty{8.4}{km}$；\qty{0.70}{}、\qty{0.35}{}、\qty{0.30}{bar}（实际为
\qty{0.26}{}：空气随高度降温）；氦为 \qty{61}{km}；火星为
\qty{10.7}{km}.
\end{solution}

\begin{solution}{exo:b2:boltzmann-factor:2}
氢：$4\eu^{-395} \approx 10^{-171}$；$4\eu^{-19.7} = 10^{-8}$；$4\eu^{-11.8} = 3 \times
10^{-5}$。钠：$\eu^{-9.8} = 6 \times 10^{-5}$。转动：$\eu^{-0.04} = 0.96$，几乎
相等。核自旋：$1 - 4 \times 10^{-6}$——磁共振所依赖的
微小极化。
\end{solution}

\begin{solution}{exo:b2:boltzmann-factor:3}
N$_2$ \qty{300}{K}: \qty{422}{}, \qty{476}{}, \qty{517}{m/s}; \qty{77}{K}: \qty{214}{}, \qty{241}{},
\qty{262}{m/s}；H$_2$：\qty{1580}{}、\qty{1780}{}、\qty{1930}{m/s}；烟粒子：
\qty{9}{cm/s}（其布朗颤动）。速率超过 $2v_{\text{p}}$ 的分子约占
二十分之一；超过 $3v_{\text{p}}$ 的约为 $2500$ 分之一——每立方米仍有 $10^{22}$ 个。
\end{solution}

\begin{solution}{exo:b2:boltzmann-factor:4}
$x = 2.2 \times 10^{-3}$ 和 $0.67$；$M/M_{\text{饱和}} = 2.2 \times 10^{-3}$ 和 $0.59$；
$M_{\text{饱和}} = n\mu = \qty{9.3e4}{A/m}$; $\chi = \mu_0n\mu^2/k_BT = 2.6 \times 10^{-4}$.
铁的 \qty{1.7e6}{A/m} 是协同效应，而非玻尔兹曼效应。
\end{solution}

\begin{solution}{exo:b2:boltzmann-factor:5}
(a) \cref{prop:b2:boltzmann-factor:twolevel}。(b) $\eu^{-x} < 1$，只要 $T > 0$；
$N_2 > N_1$ 需要 $T < 0$——即粒子数反转。(c) $0$ 与
$Nk_B\ln 2$。(d) $x = 1.93$ 时为 $0.13$；在 \qty{77}{K} 下 $x = 7.5$，为 $5 \times 10^{-4}$。
\end{solution}

\begin{solution}{exo:b2:boltzmann-factor:6}
(a) $N\varepsilon/(\eu^x + 1)$ 的导数。(b) $x = 2.40$，$C = 0.44\,Nk_B$。
(c) $\propto x^2\eu^{-x} \to 0$：冻结；$\propto 1/T^2 \to 0$：两能级均已填满，
无可吸收。(d) $k_BT = 0.42\varepsilon$：\qty{0.5}{mK}；低于
一毫开尔文时，核自旋承载了大部分热容，并
使每一次降温都变得缓慢。
\end{solution}

\begin{solution}{exo:b2:boltzmann-factor:7}
(a) \cref{thm:b2:boltzmann-factor:equipartition}。(b) $\tfrac32R$，
$\tfrac52R$。(c) $C_{\text{振动}}/R = x^2\eu^x/(\eu^x - 1)^2$：$0.002$、$0.41$、$0.90$；
$C_V = 2.50R$、$2.91R$、$3.40R$。(d) $0.92$、$0.30$、$0.995$；金刚石 $x = 4.3$：
$0.26 \times 3R = \qty{6.5}{J/mol/K}$.
\end{solution}

\begin{solution}{exo:b2:boltzmann-factor:8}
(a)、(b) \cref{thm:b2:boltzmann-factor:maxwell}。(c) $\tfrac32k_BT$。(d)
$\tfrac14n\langle v\rangle = \qty{2.9e27}{m^{-2}.s^{-1}}$：每平方
厘米每秒 $3 \times 10^{23}$ 个；通过 \qty{7.9e-13}{m^2} 的小孔：$2.3 \times 10^{15}$
个分子每秒，即泄漏率为 $\qty{1e-5}{Pa.m^3/s}$。
\end{solution}

\begin{solution}{exo:b2:boltzmann-factor:9}
(a) 对 $f$ 分部积分；主项为 $(2/\sqrt\pi)x\eu^{-x^2}$。(b)
氦 $x = 5.5$：$5 \times 10^{-13}$；氮 $x = 14.5$：$\eu^{-211}$。(c) $x =
4.9$：每碰撞时间 $10^{-10}$——已消失；$10^{-17}$ 需 $x \approx 6.6$，
$T \lesssim \qty{230}{K}$。(d) $x = 11$：$\eu^{-121}$；足够冷。
\end{solution}

\begin{solution}{exo:b2:boltzmann-factor:10}
(a) $n = n_0\eu^{m\omega^2r^2/2k_BT}$。(b) $\eu^{\Delta m\omega^2r^2/2k_BT} = \eu^{0.24} =
1.27$。(c) $H = k_BT/m'g = \qty{16}{\micro m}$；在
几个高度上计数小球，可得 $k_B$，进而 $N_A = R/k_B$。(d) 离子在电荷周围服从
$\eu^{\mp e\varphi/k_BT}$；对泊松方程作线性化，
得 $\varphi'' = \varphi/\lambda_{\text{D}}^2$：势在
$\lambda_{\text{D}}$ 之外被屏蔽。
\end{solution}

\begin{solution}{exo:b2:boltzmann-factor:11}
(a) $Z = 1/(1 - \eu^{-x})$，$\langle E\rangle = h\nu/(\eu^x - 1)$。(b) $k_BT$：
能量均分；$h\nu\eu^{-h\nu/k_BT}$：冻结。(c) $u_\nu = (8\pi\nu^2/c^3)\langle E
\rangle$；瑞利--金斯给出每个模式 $k_BT$。(d) 一个模式
不能被小于一个量子的能量激发。
\end{solution}

\begin{solution}{exo:b2:boltzmann-factor:12}
(a) 仅 $x = \mu B/k_BT$；$\ln 2$ 与 $0$。(b) $x = 0.67$ 时 $S/Nk_B = 0.51$，故
每移除一个自旋约 $0.18\,k_B$。(c) $B/T$ 不变：\qty{10}{mK}。(d) 自旋自身的
场（毫特斯拉）在末段取代 $B$，并限制降温；
测量 $M = n\mu\tanh(\mu B/k_BT)$ 即可得 $T$。
\end{solution}

\begin{solution}{pb:b2:boltzmann-factor:1}
\textbf{1.} $\dd P/\dd z = -nmg = -(mg/k_BT)P$; $H = \qty{8.4}{km}$.

\textbf{2.} $\eu^{-mgz/k_BT}$：一个分子的势能，
空气的温度。

\textbf{3.} \qty{0.70}{}, \qty{0.52}{}, \qty{0.35}{bar}; $H\ln 2 = \qty{5.8}{km}$.

\textbf{4.} $P_0/g = \qty{1.03e4}{kg/m^2}$; $\times 4\pi R^2$: \qty{5.3e18}{kg}; $\rho_0H =
1.22 \times 8400 = \qty{1.03e4}{kg/m^2}$.

\textbf{5.} $5.3 \times 10^{18}/(29 \times 10^{-3}/N_A) = 1.1 \times 10^{44}$.

\textbf{6.} $H_{\text{O}_2} = 7.6$，$H_{\text{N}_2} = \qty{8.7}{km}$：在 \qty{50}{km} 处的比值
将是其地面值的 $\eu^{-50/7.6 + 50/8.7} = 0.44$ 倍；湍流
混合的速度快于扩散的分离。

\textbf{7.} 更低：冷空气更稠密，压强下降得更快。

\textbf{8.} \qty{2.5e25}{m^{-3}}；$\eu^{-11.9}$：\qty{1.7e20}{m^{-3}}——无明显边界，只是
一种约定。

\textbf{9.} $11.2\sqrt{6370/6870} = \qty{10.8}{km/s}$.

\textbf{10.} 氦 \qty{2040}{}，\qty{2300}{m/s}；氮 \qty{770}{}，\qty{870}{m/s}。

\textbf{11.} $x = 5.3$：$4 \times 10^{-12}$；氮为 $\eu^{-196}$：无。

\textbf{12.} $\tfrac14 \times 10^{12} \times 2300 \times 4 \times 10^{-12} \times \tfrac12 \approx
10^3$ 个原子每平方米每秒；在整个地球上、经过一年，
$2 \times 10^{25}$ 个原子——\qty{0.1}{kg}。

\textbf{13.} $5 \times 10^{-6} \times 1.1 \times 10^{44} = 5.5 \times 10^{38}$ 个原子，\qty{3.7e12}{kg}；
若与 \qty{0.1}{kg/yr} 相比，滞留时间将荒谬至极。

\textbf{14.} 稳态要求损失率为 \qty{3e6}{kg/yr}：滞留
时间约为一百万年，这是公认的数值——因此
在 \qty{1000}{K} 下的热估计值远远偏小：
逃逸主要由更热的时段和非热（离子）过程主导。

\textbf{15.} $x = 3.75$：$\eu^{-14}$ 而非 $\eu^{-28}$——相差 $10^6$ 倍：
逃逸损失对外逸层温度极其敏感。

\textbf{16.} 月球既热又轻：一切气体都会在
地质年代内逃逸；土卫六温度低，对氮而言 $x = 11$。

\textbf{17.} $p_\pm = \eu^{\pm x}/2\cosh x$; $M = n\mu\tanh x$.

\textbf{18.} $x = \mu B/k_BT$：\qty{4.2}{A/m}、\qty{310}{A/m}、\qty{1.6e4}{A/m}
（饱和值 \qty{1.9e4}{A/m}）；$x \ll 1$ 时服从居里定律，即 $T \gg \qty{70}{mK}$。

\textbf{19.} $M$ 仅取决于 $T$（前提是 $B$ 已知）；变化最陡处位于
$x \approx 1$、$T \approx \mu B/k_B = \qty{70}{mK}$ 附近；
可用线圈测量磁化率。

\textbf{20.} $\langle E\rangle = -N\mu B\tanh x$；$C = Nk_Bx^2/\cosh^2x$，其峰值出现在
$x = 1.2$ 处，即 \qty{56}{mK}。

\textbf{21.} $Nk_B\ln 2$ 与 $0$：所有自旋都与磁场方向一致。

\textbf{22.} \qty{10}{mK}；自旋自身产生的毫特斯拉量级磁场取代了 $B$，
并确定了温度下限。

\textbf{23.} $x = 1$ 对应 $\mu_{\text{N}}B/k_B = \qty{0.37}{mK}$：居里定律在低至
毫开尔文量级仍然成立，温度计可用于微开尔文量级。

\textbf{24.} $\Delta\nu/\nu = \sqrt{8k_BT\ln2/mc^2} \propto \sqrt T$：
谱线宽度可用来测量气体温度。

\textbf{25.} 大气、逃逸气体、自旋：$\eu^{-E/k_BT}$，其中
$E = mgz$、$\tfrac12mv^2$、$\mp\mu B$。
\end{solution}
